\documentclass[a4paper,10pt]{article}

\usepackage[a4paper,margin=0.65in]{geometry}
\usepackage{enumitem}
\usepackage{titlesec}
\usepackage{hyperref}
\usepackage{setspace}
\usepackage{xcolor}

\setstretch{1.05}
\pagestyle{empty}
\definecolor{sectiongray}{RGB}{80,80,80}
\definecolor{dividercolor}{RGB}{200,200,200}

\titleformat{\section}{\large\bfseries\color{sectiongray}}{}{0em}{}[\titlerule]
\titlespacing*{\section}{0pt}{8pt}{5pt}

% Bullet points: add space between items and between title and first item
\setlist[itemize]{leftmargin=1.5em,itemsep=2pt,topsep=3pt,parsep=0pt,partopsep=0pt}

\setlength{\parskip}{0pt}
\setlength{\parindent}{0pt}

% Light gray divider line between projects
\newcommand{\projectdivider}{\vspace{4pt}\textcolor{dividercolor}{\hrule height 0.4pt}\vspace{4pt}}

% Project title: bold name + tech stack on same line, small gap before bullets
\newcommand{\projecttitle}[2]{\noindent\textbf{#1} \textit{#2}\par\vspace{1pt}}

% Experience sub-heading
\newcommand{\expblock}[1]{\noindent\textbf{#1}\par\vspace{1pt}}
\newcommand{\expsubblock}[1]{\noindent\textit{#1}\par\vspace{1pt}}

\begin{document}

%-------------------- Header --------------------
\begin{center}
{\LARGE\textbf{Rajnish J}}\\[3pt]
Chennai, Tamil Nadu, India\\[1pt]
\href{mailto:rajnishofficial02@gmail.com}{rajnishofficial02@gmail.com} $|$
+91-91769-60600\\[1pt]
\href{https://github.com/Rajnish-J}{GitHub} $|$
\href{https://www.linkedin.com/in/rajnish7102/}{LinkedIn} $|$
\href{https://rajnish7102.netlify.app/}{Portfolio}
\end{center}

%-------------------- Professional Summary --------------------
\section*{Professional Summary}
Software Engineer specializing in AI-powered enterprise platforms, developer productivity, infrastructure automation, and intelligent workflow orchestration. Experienced in designing scalable full-stack applications using Python, FastAPI, Next.js, React, Azure AI, OpenTofu/Terraform, LLMs, RAG, and multi-agent architectures. Strong background in enterprise platform engineering, AI-assisted software development, RBAC, system design, and developer experience.

%-------------------- Technical Skills --------------------
\section*{Technical Skills}
\textbf{Languages:} Java, Python, TypeScript, JavaScript, Go\\
\textbf{Frontend:} React.js, Next.js, Tailwind CSS\\
\textbf{Backend:} FastAPI, Spring Boot, REST APIs, WebSockets, Microservices\\
\textbf{AI:} Azure AI Foundry, Azure OpenAI, LangChain, LangGraph, RAG, MCP, AI SDK, Claude Code, GitHub Copilot, Prompt Engineering, AI Agents\\
\textbf{Infrastructure:} OpenTofu, Terraform, GitHub Actions\\
\textbf{Database:} PostgreSQL, MySQL, NeonDB\\
\textbf{Architecture:} RBAC, Authentication, Authorization, Caching, Multi-Tenant Systems, System Design

%============================================================
% PROFESSIONAL EXPERIENCE — RESTRUCTURED (PRODUCT / PROJECT / STAR)
% Assumes \expblock{}, \expsubblock{}, \projectdivider are already
% defined in your preamble (as in the original resume).
%============================================================
\section*{Professional Experience}

\noindent\textbf{Synergech Technologies, Chennai}\\
\textbf{Software Engineer} \hfill May 2025 -- Present

%------------------------------------------------------------
% PRODUCT 1: INFRAGENIE SUITE
%------------------------------------------------------------
\vspace{6pt}
\noindent\textbf{Product: \href{https://infragenie.synergech.com/auth/sign-in}{InfraGenie} Suite — Enterprise Cloud Governance \& Automation Platform}\\
\textit{Stack: Next.js, PostgreSQL, FastAPI, AI-SDK, Claude Code, OpenTofu/Terraform, Azure, ServiceNow}

\vspace{4pt}
\expblock{InfraGenie (Internal Enterprise Platform)}
\begin{itemize}
\item \textbf{Situation:} Repository provisioning, access approvals, and cloud cost oversight were handled manually across the organization, slowing DevOps delivery and creating governance gaps across Prod/Dev/QA environments.
\item \textbf{Task:} Design a single platform unifying provisioning, access governance, ticketing, and cost control for four distinct roles — platform team, end users, management, and DevOps engineers.
\item \textbf{Action:} Built an in-house RBAC \& ABAC engine, an access-package request and ticketing system, a prompt-driven OpenTofu/Terraform code-generation engine that lets DevOps engineers describe a requirement in natural language instead of consulting ChatGPT or internal documentation and have infrastructure generated and deployed directly to Azure, ServiceNow integration, environment-scoped workspaces, Intune policy management, and cost-management dashboards enforcing per-resource-group budget thresholds, using FastAPI, Next.js, and PostgreSQL.
\item \textbf{Reward:} Cut organizational cloud/billing spend by \textbf{$\sim$70\%}, reduced access-provisioning turnaround time by \textbf{$\sim$80\%}, and automated CI/CD pipeline governance for all onboarded teams.
\end{itemize}

\projectdivider

\expblock{\href{https://infra0.ai/}{infra0} (Public SaaS Platform)}
\begin{itemize}
\item \textbf{Situation:} No public-facing, lightweight equivalent of InfraGenie existed for external users who needed simplified infrastructure provisioning without full enterprise overhead.
\item \textbf{Task:} Ship a public SaaS beta carrying forward core InfraGenie capabilities — chat-driven project workflows and RBAC/ABAC — with certain enterprise-only features restricted, through a phased, version-by-version rollout.
\item \textbf{Action:} Re-architected the provisioning workflow and access-control layer into a leaner public product, gating enterprise-specific capabilities (e.g. ServiceNow integration, org-wide cost governance) behind the internal tier and iterating public feature parity release by release.
\item \textbf{Reward:} Delivered a functioning public beta with a clear version-wise feature-parity roadmap, extending platform reach beyond internal enterprise users.
\end{itemize}

\projectdivider

\expblock{Infra Landing Platform}
\begin{itemize}
\item \textbf{Situation:} infra0 had no external-facing channel to communicate its features, pricing, or workflows to prospective users.
\item \textbf{Task:} Build a marketing landing site that converts visitors into signups by clearly showcasing product value.
\item \textbf{Action:} Designed and developed a responsive landing website covering feature highlights, pricing tiers, and workflow walkthroughs; produced supporting LinkedIn marketing content and a recorded product demo video for go-to-market.
\item \textbf{Reward:} Established the primary acquisition and marketing entry point for the public product launch.
\end{itemize}

\projectdivider

\expblock{Infra Documentation Platform}
\begin{itemize}
\item \textbf{Situation:} Public users lacked a self-service resource to learn or troubleshoot the platform, increasing support dependency.
\item \textbf{Task:} Deliver a documentation portal supporting search, multilingual access, and AI-assisted content help.
\item \textbf{Action:} Built a documentation portal with Markdown rendering, full-text search, multilingual support, and AI-assisted documentation generation.
\item \textbf{Reward:} Improved user self-service and onboarding, reducing inbound support load for the public product.
\end{itemize}

%------------------------------------------------------------
% PRODUCT 2: AGENTIC WORKFLOW ORCHESTRATION & ANALYTICS PLATFORM
%------------------------------------------------------------
\vspace{10pt}
\noindent\textbf{Product: Agent Orchestration \& Analytics Platform (Agents Marketplace)}\\
\textit{Stack: Next.js, Python, PostgreSQL, VS Code Extension API, Claude Code, Azure AI Foundry}

\vspace{4pt}
\expblock{Claude Code Analytics \& Marketplace Platform}
\begin{itemize}
\item \textbf{Situation:} Leadership had no visibility into Claude Code, GitHub Copilot, and Codex seat usage or developer productivity, and teams had no centralized way to reuse AI skills/agents across the organization.
\item \textbf{Task:} Build an analytics and governance layer giving management usage insight while enabling org-wide, version-controlled agent/skill reuse.
\item \textbf{Action:} Built a usage-upload and insights engine, a Skills \& Agents Marketplace with fork mechanism and version control, seat-allocation management, live and historical dashboards, org-wide RBAC, and Azure AI Foundry agent-registry integration.
\item \textbf{Reward:} Gave leadership real-time visibility into AI-tool adoption and developer productivity, enabling data-driven seat reallocation across \textbf{100\%} of onboarded teams.
\end{itemize}

\vspace{4pt}
\expsubblock{Workflow Orchestration (Business Users)}
\begin{itemize}
\item \textbf{Situation:} Business teams had no no-code way to chain Azure AI Foundry agents into repeatable, multi-step workflows.
\item \textbf{Task:} Build an orchestration engine letting non-developers create, pull from marketplace, and run agent-based workflows.
\item \textbf{Action:} Implemented a no-code workflow builder, execution engine, monitoring, and orchestration pipelines connecting marketplace and Foundry agents; work is now extending into agent evaluation and governance.
\item \textbf{Reward:} Enabled business users to automate multi-step tasks end-to-end without engineering support.
\end{itemize}

\vspace{4pt}
\expsubblock{VS Code Extension (Developer SDLC)}
\begin{itemize}
\item \textbf{Situation:} Developers lacked an equivalent agent-orchestration capability embedded in their IDE for build–develop–deploy workflows.
\item \textbf{Task:} Deliver a VS Code extension enabling in-IDE creation and execution of AI-driven SDLC workflows.
\item \textbf{Action:} Built an extension supporting agent install/pull, configurable workflow-driven code generation, and orchestrator-triggered automation across the SDLC.
\item \textbf{Reward:} Reduced developer context-switching and standardized AI-assisted automation directly inside the IDE.
\end{itemize}

\projectdivider

\expblock{Autonomous Testing Platform}
\begin{itemize}
\item \textbf{Situation:} QA work was split across manual task tracking and separate API-testing tools, with no collaborative, AI-assisted testing workspace.
\item \textbf{Task:} Build a unified workspace combining project/task management with autonomous, AI-driven test execution.
\item \textbf{Action:} Developed collaborative Kanban-based workspaces and integrated Playwright MCP, Postman MCP, and AI chat agents for autonomous UI/API test authoring and execution.
\item \textbf{Reward:} Consolidated test planning and execution into a single AI-assisted workspace, cutting manual test-authoring effort.
\end{itemize}

%------------------------------------------------------------
% PRODUCT 3: RISKGUARD AI (CLIENT PROJECT)
%------------------------------------------------------------
\vspace{10pt}
\noindent\textbf{Project: RiskGuard AI (Client Project) — AI-Assisted Underwriting Platform}\\
\textit{Stack: Python (FastAPI), React.js, PostgreSQL, Azure AI Foundry, Azure AI Search, Azure Content Understanding, Claude Code, GitHub Copilot}

\vspace{4pt}
\expblock{Document-Driven Underwriting Automation}
\begin{itemize}
\item \textbf{Situation:} Underwriters manually reviewed applicant documents and researched risk context before approving or denying insurance claims, taking anywhere from half a day to over two days per case.
\item \textbf{Task:} Build a human-in-the-loop (HITL) system that automates document verification, contextual research, and decision support while keeping underwriters in control of the final call.
\item \textbf{Action:} Built a document-processing pipeline on Azure AI Search and Azure Content Understanding, integrated automated web research for supporting insights, and developed a React-based workspace with multi-document project versioning for underwriters.
\item \textbf{Reward:} Reduced typical review time from up to \textbf{2 days} to as little as \textbf{$\sim$15 minutes}, with edge cases capped at $\sim$1 hour — a time reduction of roughly \textbf{90--98\%}.
\end{itemize}

%------------------------------------------------------------
% EARLIER ROLES
%------------------------------------------------------------
\vspace{10pt}
\noindent\textbf{Software Engineer Intern} \hfill Feb 2025 -- Apr 2025
\begin{itemize}
\item Built AI-assisted document processing, workflow automation, and backend integrations.
\item Worked on enterprise proof-of-concepts involving AI, APIs, and scalable applications.
\end{itemize}

\vspace{4pt}
\noindent\textbf{Graduate Engineer Trainee} \hfill Aug 2024 -- Jan 2025
\begin{itemize}
\item Developed Spring Boot and React applications with REST APIs.
\item Strengthened backend, database, and system architecture fundamentals.
\end{itemize}

%------------------------------------------------------------
\section*{AI-Assisted Software Engineering}
\begin{itemize}
\item Advanced user of Claude Code and GitHub Copilot for enterprise software development.
\item Experienced with Agent Mode, Edit Mode, Planning Mode, approval workflows, and AI-driven development practices.
\item Integrated reusable Claude Skills and AI Agents into engineering workflows.
\item Built AI-assisted workflows for development, documentation, testing, Git operations, commits, and pull requests.
\item Selected appropriate AI models based on reasoning complexity, speed, and context requirements.
\item Experienced with MCP integrations, AI SDKs, RAG pipelines, and enterprise agent workflows.
\end{itemize}

%-------------------- Leadership --------------------
\section*{Leadership \& Technical Contributions}
\begin{itemize}
\item Led Next.js architecture migration and established reusable project structure.
\item Conducted internal knowledge sharing on Next.js, AI SDKs and AI development practices.
\item Participated in client demos, solution presentations and architecture discussions.
\item Contributed to coding standards, reusable architecture and developer productivity initiatives.
\end{itemize}

%-------------------- Personal Projects --------------------
\section*{Personal Projects}

\projecttitle{Harness --- AI Coding-Agent Platform}{--- Python, FastAPI, TypeScript, Next.js, PostgreSQL, Docker, LangGraph}

Designed and independently built Harness, a full-stack AI coding-agent platform, to understand production agent systems end-to-end. The architecture cleanly separates a Python/FastAPI backend --- owning the agent's decide $\rightarrow$ act $\rightarrow$ observe loop, tool execution, and safety guardrails --- from a Next.js/TypeScript frontend, communicating over a streamed HTTP boundary.

Iteratively extended the system from a single-tool chat prototype into a multi-surface platform: a provider-agnostic LLM abstraction (Anthropic/OpenAI) that keeps each provider's native tool-calling shape rather than a translated common format; a sandboxed tool layer for file, git, shell, and search operations; a visual workflow builder for composing and running multi-agent pipelines (LangGraph DAG engine plus a React Flow canvas, checkpointed to PostgreSQL via Drizzle ORM); Model Context Protocol support for external tool servers; and an AES-256-GCM encrypted credentials vault for third-party integrations.

Currently building project-level isolation: the platform clones a user's GitHub repository into a managed workspace, indexes its files, and provisions a dedicated Docker container per project so agent-run shell commands (installs, builds, tests) execute in a reproducible sandbox --- with transparent, automatic fallback to host execution when Docker is unavailable.

Built to learn by doing: several AI coding-agent harnesses already exist in the market, but building one from scratch --- rather than only using existing ones --- was the most direct way to understand how large AI platforms actually engineer these systems under the hood.

\projectdivider

\projecttitle{\href{https://task-tracker-gray-mu.vercel.app/}{Task Tracker} --- Full-Stack AI-Augmented Project Management Platform}{--- Next.js 16 (App Router), React 19, TypeScript, PostgreSQL (Neon, serverless), Drizzle ORM, better-auth, Anthropic Claude API, Tailwind CSS, dnd-kit, TanStack Table, Recharts}

An end-to-end task and project management web application, similar in spirit to Trello/Linear, built to explore production-grade patterns in full-stack architecture, access control, and LLM-powered application design. Built first to solve a personal problem: struggling to track personal progress, in-flight work, and workload across projects, so an analytical dashboard and timeline view to see it all at a glance felt worth building --- reflecting a broader preference for building tools that solve a real problem for the user and the community rather than paying for an equivalent commercial product.

Key contributions:
\begin{itemize}
\item Modeled a normalized data hierarchy (Section $\rightarrow$ Project $\rightarrow$ Column $\rightarrow$ Task $\rightarrow$ Checklist Item) in PostgreSQL via Drizzle ORM, with cascade/set-null delete rules, unique partial indexes, and safe migration history (4 schema migrations to date).
\item Built a multi-tenant workspace model supporting both personal accounts and shared team spaces, with role-based membership, pending/accepted/declined invitations, and an in-app notification system.
\item Designed a fine-grained permission engine scoping create/update/delete rights independently across projects, sections, columns, and tasks --- including per-project grants versus team-wide defaults.
\item Engineered an embedded AI chat assistant on Anthropic's Claude API using structured tool-use (18+ tools across read/write operations), prompt caching to separate a stable system prompt from a volatile workspace snapshot, and explicit guardrails that keep the agent scoped to the app's domain and resistant to prompt injection from user-entered task content.
\item Implemented authentication and session handling with better-auth, including cookie-based route protection middleware and secure email/password flows.
\item Built interactive UI features: drag-and-drop Kanban boards (dnd-kit), a tag-filterable timeline view, an analytics dashboard (Recharts), batch operations for checklist items, and a custom shadcn/ui-based component/theming system.
\end{itemize}

\noindent\textit{Currently working on:} refining frontend polish and visual identity (custom WebGL background effects, redesigned login experience); consolidating and refactoring the codebase for long-term readability as the feature surface grows; and extending user-configurable preferences (task defaults, week-start day) with reusable UI patterns.

\projectdivider

\projecttitle{useApiGuard}{--- Next.js, TypeScript, React}
\begin{itemize}
\item Designed and published a reusable React hook to standardize API communication across applications.
\item Simplified API handling by centralizing loading states, error handling, request lifecycle management, and response processing.
\item Open-sourced the package with comprehensive documentation and real-world implementation examples.
\item Published as an NPM package for reuse across multiple React and Next.js projects.
\end{itemize}

\projectdivider

\projecttitle{Hospify -- Hospital Management System}{--- Spring Boot, React.js, MySQL}
\begin{itemize}
\item Developed a comprehensive hospital management platform supporting patients, doctors, hospitals, and administrative workflows.
\item Built modules for patient registration, doctor management, appointment scheduling, and hospital administration.
\item Implemented seamless appointment booking with doctor availability management and scheduling workflows.
\item Designed medicine management and prescription tracking modules to streamline healthcare operations.
\item Developed RESTful APIs using Spring Boot and integrated them with a responsive React.js frontend.
\item Designed relational database schemas and managed source control using Git and GitHub.
\end{itemize}

%-------------------- Achievements --------------------
\section*{Achievements}
\begin{itemize}
\item Emerging Star Award -- Synergech (2025)
\end{itemize}

%-------------------- Certifications --------------------
\section*{Certifications \& Courses}
\begin{itemize}
\item Claude Code in Action -- Anthropic
\item Building with the Claude API -- Anthropic
\item Introduction to Model Context Protocol -- Anthropic
\item Model Context Protocol: Advanced Topics -- Anthropic
\item Introduction to Agent Skills
\item Microsoft Certified: Azure AI Engineer Associate (AI-102) -- Preparing
\end{itemize}

\end{document}
